\begin{thebibliography}{99}
\raggedright

\bibitem{CLZZ}
H.~Chen, Z.~Li, R.~Zhang, and X.~Zhang,
The number of limit cycles of piecewise linear Li\'enard systems,
\href{https://arxiv.org/abs/2608.19542v1}{arXiv:2608.19542v1}
[math.DS], 20 August 2026.

\bibitem{CGP}
B.~Coll, A.~Gasull, and R.~Prohens,
Limit cycles for non smooth differential equations via Schwarzian derivative,
\emph{Journal of Differential Equations} \textbf{132}(2)
(1996), 203--221.
\href{https://doi.org/10.1006/jdeq.1996.0177}{doi:10.1006/jdeq.1996.0177}.

\bibitem{DMD}
P.~De Maesschalck and F.~Dumortier,
Classical Li\'enard equations of degree $n\ge6$ can have
$\lfloor(n-1)/2\rfloor+2$ limit cycles,
\emph{Journal of Differential Equations} \textbf{250}(4)
(2011), 2162--2176.
\href{https://doi.org/10.1016/j.jde.2010.12.003}{doi:10.1016/j.jde.2010.12.003}.

\bibitem{DMH}
P.~De Maesschalck and R.~Huzak,
Slow divergence integrals in classical Li\'enard equations near centers,
\emph{Journal of Dynamics and Differential Equations}
\textbf{27}(1) (2015), 177--185.
\href{https://doi.org/10.1007/s10884-014-9358-1}{doi:10.1007/s10884-014-9358-1}.

\bibitem{DPR}
F.~Dumortier, D.~Panazzolo, and R.~Roussarie,
More limit cycles than expected in Li\'enard equations,
\emph{Proceedings of the American Mathematical Society}
\textbf{135}(6) (2007), 1895--1904.
\href{https://doi.org/10.1090/S0002-9939-07-08688-1}{doi:10.1090/S0002-9939-07-08688-1}.

\bibitem{GGG}
A.~Gasull, H.~Giacomini, and M.~Grau,
Effective construction of Poincar\'e--Bendixson regions,
\emph{Journal of Applied Analysis and Computation}
\textbf{7}(4) (2017), 1549--1569.
\href{https://doi.org/10.11948/2017094}{doi:10.11948/2017094}.

\bibitem{Rosales}
M.~Hern\'andez Rosales,
A global algebraic variety approach to the maximum number of limit
cycles in Li\'enard equations,
preprint, 31 August 2026, Section~9, p.~32.
\href{https://www.researchgate.net/profile/Manuel-Rosales-2/publication/404693751_A_bound_for_the_maximum_number_of_limit_cycles_of_the_Lienard_equation/links/6ac3d9555770bb0daeda4d14/A-bound-for-the-maximum-number-of-limit-cycles-of-the-Lienard-equation.pdf}{author's manuscript}.

\bibitem{LiLlibre}
C.~Li and J.~Llibre,
Uniqueness of limit cycles for Li\'enard differential equations
of degree four,
\emph{Journal of Differential Equations} \textbf{252}(4)
(2012), 3142--3162.
\href{https://doi.org/10.1016/j.jde.2011.11.002}{doi:10.1016/j.jde.2011.11.002}.

\bibitem{LiLu}
C.~Li and K.~Lu,
Slow divergence integral and its application to classical
Li\'enard equations of degree 5,
\emph{Journal of Differential Equations} \textbf{257}(12)
(2014), 4437--4469.
\href{https://doi.org/10.1016/j.jde.2014.08.015}{doi:10.1016/j.jde.2014.08.015}.

\bibitem{LMP}
A.~Lins, W.~de Melo, and C.~C.~Pugh,
On Li\'enard's equation,
in \emph{Geometry and Topology}, J.~Palis and M.~do Carmo (eds.),
Lecture Notes in Mathematics \textbf{597},
Springer, Berlin, 1977, pp.~335--357.
\href{https://doi.org/10.1007/BFb0085364}{doi:10.1007/BFb0085364}.

\bibitem{LlibreZhang}
J.~Llibre and X.~Zhang,
Limit cycles of the classical Li\'enard differential systems:
A survey on the Lins Neto, de Melo and Pugh's conjecture,
\emph{Expositiones Mathematicae} \textbf{35}(3) (2017), 286--299.
\href{https://doi.org/10.1016/j.exmath.2016.12.001}{doi:10.1016/j.exmath.2016.12.001}.

\bibitem{Odani}
K.~Odani,
Existence of exactly $N$ periodic solutions for Li\'enard systems,
\emph{Funkcialaj Ekvacioj} \textbf{39}(2) (1996), 217--234.
\href{https://doi.org/10.24546/0100499900}{doi:10.24546/0100499900}.

\bibitem{OT}
V.~Ovsienko and S.~Tabachnikov,
What Is \ldots\ the Schwarzian Derivative?,
\emph{Notices of the American Mathematical Society}
\textbf{56}(1) (2009), 34--36.
\href{https://ovsienko.perso.math.cnrs.fr/Publis/what-is.pdf}{Author's text}.

\bibitem{Perko}
L.~Perko,
\emph{Differential Equations and Dynamical Systems},
3rd ed., Texts in Applied Mathematics \textbf{7},
Springer, New York, 2001.
\href{https://doi.org/10.1007/978-1-4613-0003-8}{doi:10.1007/978-1-4613-0003-8}.

\bibitem{Rychkov}
G.~S.~Rychkov,
The maximal number of limit cycles of the system
$\dot y=-x$, $\dot x=y-\sum_{i=0}^{2}a_i x^{2i+1}$
is equal to two,
\emph{Differentsial'nye Uravneniya} \textbf{11}(2) (1975),
390--391 (in Russian).
\href{https://www.mathnet.ru/eng/de2400}{Math-Net: de2400}.

\end{thebibliography}
